% ===== GAN artifact table (N_train=10000, N_eval=2000) =====
\begin{table}[t]
  \centering
  \begin{tabular}{lrrrr}
    \toprule
    \textbf{Method} & Follow-CNN & Anti-CNN & Reported & $F_{\text{rank}}$ \\
     & $\mathrm{ASR}@5$ & $\mathrm{ASR}@5$ & $\mathrm{ASR}@5$ & \\
    \midrule
    HoneyGen Baseline & 0.746 & 0.038 & 0.746 & 0.169 \\
    HoneyGen + Prescreen-Ratio & 0.341 & 0.180 & 0.341 & 0.422 \\
    HoneyGen + Hybrid-Ratio & 0.293 & 0.202 & 0.293 & 0.462 \\
    GAN (Self-Trained) & 0.025 & 0.872 & 0.872 & 0.920 \\
    \bottomrule
  \end{tabular}
  \caption{Two attacker strategies at $N_{\text{train}}=10000$:
  \emph{Follow-CNN} guesses the CNN's top-ranked candidates,
  \emph{Anti-CNN} guesses its bottom-ranked ones, and \emph{Reported}
  is the rational-attacker value of \Cref{eq:rational}. For the baseline
  and the ratio filters, Follow-CNN is the larger of the two directions
  and $F_{\text{rank}}$ sits near or above $0.4$, so the reported value
  simply follows the CNN. The GAN is the exception --- its
  $F_{\text{rank}}$ near $1$ places the real password at the bottom, so
  Follow-CNN succeeds only $2.5\%$ of the time while Anti-CNN reaches
  $87\%$. Its low Follow-CNN score is inverted distinguishability, not
  flatness.}
  \label{tab:gan}
\end{table}
